\section{Use Case Specifications}

\newcommand{\UseCase}[8]{%
\begin{center}
\renewcommand{\arraystretch}{1.25}
\begin{tabular}{|p{0.95\textwidth}|}
\hline
\centering\textbf{Прецедент: #2} \tabularnewline
\hline
\textbf{ID:} #1 \tabularnewline
\hline
\textbf{Краткое описание:} #3 \tabularnewline
\hline
\textbf{Главный актер:} #4 \tabularnewline
\hline
\textbf{Предусловия:} #5 \tabularnewline
\hline
\textbf{Основной поток:}\newline #6 \tabularnewline
\hline
\textbf{Альтернативный поток:} #7 \tabularnewline
\hline
\textbf{Постусловия:} #8 \tabularnewline
\hline
\end{tabular}
\end{center}
\vspace{0.4cm}
}

\subsection{Guest}

\UseCase
{1}
{Просмотр календаря}
{Гость просматривает календарную сетку и выбирает интересующую дату}
{Гость}
{Открыта главная страница портала}
{1. Система отображает текущий месяц.\newline
2. Гость листает месяцы и выбирает дату.\newline
3. Система показывает данные по выбранной дате.}
{Календарные данные недоступны: система выводит сообщение об ошибке и предлагает повторить попытку.}
{Гость получил актуальную календарную информацию по выбранной дате.}

\UseCase
{2}
{Просмотр праздников}
{Гость просматривает список праздников на выбранную дату}
{Гость}
{Выбрана дата в календаре}
{1. Гость переходит в раздел праздников.\newline
2. Система отображает список праздников и их типы.\newline
3. Гость открывает карточку праздника для деталей.}
{Для выбранной даты праздники отсутствуют: система показывает пустой список.}
{Гость просмотрел список и описание праздников.}

\UseCase
{4}
{Просмотр гороскопов}
{Гость просматривает гороскопы по знакам зодиака}
{Гость}
{Открыт раздел гороскопов}
{1. Гость выбирает знак зодиака.\newline
2. Система отображает прогноз на выбранный период.\newline
3. Гость просматривает прогноз по категориям.}
{Контент временно недоступен: система предлагает повторить запрос позже.}
{Гость просмотрел гороскоп для выбранного знака.}

\UseCase
{5}
{Просмотр исторических событий}
{Гость просматривает исторические события, связанные с датой}
{Гость}
{Открыт раздел исторических событий}
{1. Гость выбирает дату.\newline
2. Система выводит список событий по дате.\newline
3. Гость открывает подробности события.}
{События по выбранной дате отсутствуют: система показывает пустой результат.}
{Гость ознакомился с историческими событиями.}

\UseCase
{6}
{Поиск контента}
{Гость выполняет поиск по ключевым словам, датам и категориям}
{Гость}
{Открыта страница с поисковой строкой}
{1. Гость вводит запрос и при необходимости выбирает фильтры.\newline
2. Система выполняет поиск.\newline
3. Система показывает результаты с типами контента.}
{Ничего не найдено: система предлагает скорректировать запрос.}
{Гость получил релевантный список материалов.}
